\section{Scope, Method, and Qualifications}

\subsection{Reader, question, and claim}

This companion essay to \emph{The Creature Before God} is a discursive philosophical and theological article for a serious general reader. It asks whether the success of physical instruments warrants the conclusion that only instrument-detectable realities are real or that only instrument-confirmable propositions are knowledge. It argues that it does not: indications become measurement results only within rationally governed practices of definition, calibration, modelling, inference, and judgment, so the success of measurement cannot establish measurement as the exhaustive criterion of reality or truth.

The article uses \emph{instrumentalist} in a deliberately coined and narrow sense. \emph{Instrumental exclusivism} names the conversion of an appropriate scientific method into an exclusive epistemology or ontology. It is not the established philosophy-of-science position called instrumentalism, which concerns the status and function of scientific theories; it is not ordinary empiricism; and it is not identical with metaphysical naturalism. Scientism, verificationism, and operationalism are treated as related but distinct positions. The coined term and the organization of the argument are project synthesis, not received technical definitions or an ecclesial classification.

\subsection{Sources and argumentative roles}

\begin{threestable}{Source class}{Function in the article}{Governing boundary}
Conciliar and papal teaching & Defends the proper autonomy of science, rejects scientism, and affirms the reach of reason from empirical phenomena toward metaphysical foundation. & These sources govern the Catholic intellectual boundary; they do not turn a particular philosophy of measurement into doctrine.\\
Thomistic philosophy and theology & Supplies the account of knowledge beginning in sense, intellectual abstraction, and natural knowledge of God from effects without comprehension of the divine essence. & The adopted Thomistic articulation is distinguished from dogma and is not represented as the only permitted Catholic school.\\
Philosophy and practice of science & Distinguishes indications, measurement results, calibration, operationalism, instrumentalism, and verificationism. & Technical and reference sources describe concepts and debates; they neither prove theology nor receive magisterial authority.\\
Modern apologetic witness & Supplies Lewis's cause--ground pressure and the history of its revision after Anscombe's criticism. & The article uses a narrow analogy and does not claim that Lewis refuted every naturalism or demonstrated an immaterial mind.\\
Project synthesis & Coins instrumental exclusivism, applies the cause--ground distinction to measurement, and orders the movement from method to metaphysics. & The synthesis clears an epistemic obstacle; it does not demonstrate God or authenticate a miracle.\\
\end{threestable}

\subsection{Included and excluded scope}

Included are the legitimate scope of instrument-mediated inquiry; the distinction between instrument indication and measurement result; calibration, models, inference, and uncertainty as conditions of measurement; the self-application and domain-success objections to an exclusive criterion; the difference among several modes of rational warrant; the Thomistic movement from sensible particulars to intellectual knowledge; and the distinction between empirical causes and metaphysical inquiry into foundation.

Excluded are a comprehensive history of logical positivism or philosophy of science; adjudication between scientific realism and antirealism; a complete argument from reason; a philosophy of mind, consciousness, mathematics, or moral epistemology; arguments from quantum mechanics; a defense of any reported miracle; and the demonstrations for God's existence developed elsewhere in the series. The article does not claim that every unmeasurable assertion is meaningful or true, that scientific methods are unreliable, or that empirical evidence is irrelevant to historical, moral, philosophical, or theological judgment. It makes no legal, jurisdictional, or mutable disciplinary conclusion.

\subsection{Material qualifications}

\begin{enumerate}[leftmargin=1.45em,itemsep=0.25em]
\item Instrument indications are not dismissed as ``mere data.'' Their evidential force is explained through the practices that establish what they indicate.
\item Theory-dependence does not entail that measurements are subjective or circular. Models and calibration make results criticizable, comparable, and corrigible.
\item The criterion's failure to satisfy its own demand is a real pressure, but the article also answers the stronger version in which instrumental exclusivism is proposed as a convention justified by scientific success.
\item ``Measurable in principle'' is not assumed to mean direct visibility. Instruments often warrant claims about unobservable entities through effects and models.
\item Rejecting an exclusive instrumental criterion does not validate superstition, private revelation, intuition, or metaphysical assertion without evidence. Claims remain accountable to warrants fitted to their objects.
\item The moral example establishes that technical power does not generate its own normative limit. It does not deny that empirical facts are indispensable to sound moral judgment.
\item Natural reason's movement from sensible effects to God is adopted from Aquinas. The article clears the method's admissibility but neither reproduces the demonstrations nor treats their conclusions as instrument readings.
\item God is not placed in a gap left by incomplete science. The adopted classical claim concerns the cause of the whole created order and therefore does not compete with secondary explanations within that order.
\end{enumerate}

\subsection{Rights, currentness, and review}

The online official, technical, and scholarly witnesses were checked on 23 July 2026. No Lewis wording is quoted. The account of his argument and its revision is a focused paraphrase supported by the C.~S. Lewis Foundation study guide and the identified Oxford records; the precise wording of a licensed modern edition was not independently collated for this draft. The brief phrase ``from phenomenon to foundation'' is attributed to \emph{Fides et ratio}; official text, third-party wording, and translations retain their own status and are not relicensed by the project's CC BY 4.0 grant. The surrounding prose and synthesis are project-created.

This is an internally source-audited working study. Production and artifact evidence for the current source is recorded in the repository research scope. The installed snapshot remains on hold pending exact-snapshot authorization. Independent philosophy-of-science, Thomistic, theological, literary, scientific, and ecclesiastical review remain outstanding; no release clearance, imprimatur, nihil obstat, or ecclesiastical approval is claimed.
